\section{Marked periods of a highest Hodge vector}
\label{sec:marked-periods}

The vector supplied by the covering construction need not generate the
highest Hodge space.  We therefore retain the whole highest space and add
flags determined by the vector on certain compact factors.  This section
constructs the quotient and its line bundles.  The boundary and Higgs-tail
arguments in \Cref{sec:boundary-rank} complete the positivity statement.

\begin{proposition}[Marked-period construction]\label{prop:marked-period}
Let $B^\circ$ be a smooth irreducible complex quasi-projective variety,
and let $\mathbb V$ be a polarizable rational pure variation with an
invariant integral lattice.  Suppose that a line bundle $\mathcal L$ has
a generically nonzero algebraic morphism
\[
 u:\mathcal L\longrightarrow F^p\mathbb V_{\C},
 \qquad F^{p+1}\mathbb V_{\C}=0.
\]
There are a positive integer $a$, a homogeneous tensor construction
followed by flat Hodge-compatible projections, and a generically nonzero
morphism
\begin{equation}\label{eq:marked-vector}
 u_*:\mathcal L^{\otimes a}\longrightarrow E_0,
 \qquad E_0=F^{p_*}\mathbb A,
 \qquad F^{p_*+1}\mathbb A=0,
\end{equation}
where $\mathbb A$ is a polarized complex pure variation.  The operations
may first be performed equivariantly at finite monodromy level.  They
descend to $B^\circ$ after shrinking it.  On smooth projective
birational models there is a dominant morphism
$b:B\to S$ with connected fibers and the following properties.
\begin{enumerate}
\item \emph{Descent.} On dense smooth opens, $\mathbb A$ is the pullback of an actual
variation on $S^\circ$.  Its connected representation is irreducible
and factors through a rational semisimple adjoint group $G$.  The full
monodromy acts through its projected action on
$\mathfrak g=\Lie G$.  In particular, finite scalar monodromy in the
kernel of that action has been eliminated.
\item \emph{The marks.} There is an algebraic compact-flag section $j_{\mathrm u}$,
possibly empty, determined pointwise and equivariantly by $[u_*]$.
In flat coordinates the joint marks
\[
 j=(E_0,j_{\mathrm u})
\]
take values in a closed product flag orbit $J$ of $G_{\C}$.  The two
kinds of marks use disjoint complex simple factors.  Every rational
simple factor of $G$ has a complex factor acting nontrivially on $J$.
The full Lie-period map together with $j_{\mathrm u}$ is generically
immersive on $S^\circ$.
\item \emph{The nef lines and adjoint positivity.} If $S$ is a point, $\mathbb A$ is constant and we set
$\mathcal H=\mathcal H_{\mathrm u}=0$.  Otherwise $S$ has
nef rational line bundles $\mathcal H$ and $\mathcal H_{\mathrm u}$,
with $\mathcal H-\mathcal H_{\mathrm u}$ nef.  On a smooth unipotent
connected-monodromy cover their pullbacks are respectively
\[
 \det E_0+\mathcal H_{\mathrm u}
 \quad\hbox{and the positive compact-flag line.}
\]
The generic curvature rank of $\mathcal H$ is the generic rank of
$dj$, and equals $\nu(\mathcal H)$.  If $D_0$ denotes the reduced
union of boundary components with nonidentity local monodromy on
$\mathfrak g$, then
\begin{equation}\label{eq:marked-initial-big}
 K_S+D_0+c\mathcal H_{\mathrm u}\quad\text{is big for }c\gg0.
\end{equation}
Furthermore,
\begin{equation}\label{eq:marked-final-big}
 K_S+b_0\mathcal H\quad\text{is big for }b_0\gg0.
\end{equation}
\end{enumerate}
All assertions about extensions are compatible with further smooth
projective modifications and pullback at unipotent level.  In
particular, the iterated Higgs fields of $u_*$ are polynomial tensor
expressions in those of $u$; the construction introduces no division
by a local function.
\end{proposition}

We prove the construction and \eqref{eq:marked-initial-big} below.
Equation \eqref{eq:marked-final-big} is
\Cref{prop:adjoint-bigness}.  The other comparison needed later concerns
a chosen line in a pullback of $E_0$: \Cref{prop:tail-rank} constructs
a nef line $P$ from its iterated Higgs images and proves that adding
the pulled-back $\mathcal H$ does not increase $\nu(P)$, provided the
chosen line induces the same compact marks.  Both comparisons are
proved in the next section, after the quotient is constructed.

\subsection{Monodromy, representations, and homogeneous replacements}

Write $M$ for the full rational algebraic monodromy group of
$\mathbb V$, and $G_c=M^\circ$.  A connected monodromy level means a
connected finite étale cover on which the monodromy lies in $G_c$;
its image is still Zariski dense there.

\begin{lemma}\label{lem:marked-monodromy}
The group $G_c$ is semisimple.  At connected monodromy level the Hodge
circles normalize $G_c$, are locally conjugate by $G_c(\mathbb R)^+$,
and have locally constant action on every $G_c$-invariant tensor
space.  The horizontal period differential therefore has the form
\begin{equation}\label{eq:marked-period-tangent}
 t(v)\in\mathfrak g_{c,\C}^{-1,1},
 \qquad \theta(v)=d\rho\bigl(t(v)\bigr),
 \qquad [t(v),t(w)]=0.
\end{equation}
The Lie algebra, and every monodromy-stable sum of its rational
simple ideals, is a polarized rational integral variation of weight
zero.  Its Hodge metric makes $x\mapsto-x^*$ a compact conjugation,
preserving each complex simple ideal.
\end{lemma}

\begin{proof}
André's monodromy normality theorem
\cite[\S5, Theorem~1]{Andre92} applies to this polarizable pure
variation over a smooth connected algebraic variety: the connected
monodromy is normal in the derived generic Mumford--Tate group.
It is therefore semisimple.  The same statement applies after a
finite étale cover and after restriction to a smooth algebraic
subvariety, with its own generic Mumford--Tate group.

The theorem of the fixed part for polarized variations says that
the invariant part of each tensor variation is a constant Hodge
structure; see \cite[Theorem~7.22 and Corollary~7.23]{Schmid}.
These tensor invariants determine
the reductive group $G_c$.  To see the latter assertion in the form
needed here, a $G_c$-stable subspace in a tensor representation has
a $G_c$-equivariant complementary subspace, so it can be encoded by
its invariant projector.  The tensor-stabilizer description of an
algebraic subgroup then applies
\cite[Theorem~9.2 and \S\S9.6--9.7]{Milne}.  Thus, in a fixed flat frame, the
Hodge circles $h_x$ normalize $G_c$, and
\[
 h_x(z)h_{x_0}(z)^{-1}\in G_c(\mathbb R).
\]
In particular their images in the quotient by $G_c$ are constant.
For completeness, local conjugacy follows directly from compactness
of the circle.  Along a smooth path of these homomorphisms, the
logarithmic derivative is a Lie-algebra-valued cocycle for the
circle action.  Averaging a cocycle over the circle writes it as a
coboundary.  Integrating the resulting time-dependent element of
$\Lie G_c(\mathbb R)$ conjugates the circles along the path.  One
can do this first in the adjoint group and lift through the central
isogeny; the remaining central discrepancy is locally constant.

The tangent to the filtration orbit is obtained from the negative
degrees in $\mathfrak g_{c,\C}$.  Griffiths transversality makes its
only possible degree equal to $-1$, giving $t(v)$ and its indicated
action in every representation.  Higgs integrability gives the
commutation relation in \eqref{eq:marked-period-tangent}.

The Lie algebra sits as a rational flat subbundle of
$\End\mathbb V$.  Normalization by the Hodge circles makes it a
Hodge subvariation, and restriction of the endomorphism polarization
polarizes it.  Intersecting its rational fiber with the endomorphism
lattice supplies an invariant full lattice.  These assertions also
hold for the stated rational sums of ideals.  Finally, preservation
of the polarization and normalization by the Hodge grading imply
stability under Hodge-metric adjoint.  On the real Lie algebra the
negative adjoint is the Cartan involution obtained from the Weil
operator.  Its complex antilinear extension is the claimed compact
conjugation.  A compact real form of a semisimple complex algebra
is the direct sum of compact real forms of its simple ideals, so
this conjugation preserves each ideal.
\end{proof}

\begin{lemma}[Homogeneous replacement]\label{lem:marked-replacement}
The replacement \eqref{eq:marked-vector} can be chosen so that its
connected representation is irreducible and its full representation
factors through the action on the retained rational adjoint Lie
factors.  Among all such homogeneous replacements, choose one with
the fewest retained rational simple factors.  Write $G$ for their
product as an adjoint group and $\mathfrak g=\Lie G$.

Every representation factor used in splitting $\mathbb A$ according
to a product decomposition of $G$ is a Hodge-compatible summand of
tensor constructions on the corresponding Lie variation, up to a
constant shift of the Hodge indices.  It can be equipped with the
metric for which representation adjoints are induced by Lie
adjoints.  The highest space $E_0$ is the tensor product of the
highest spaces of these factors and has a closed projective
homogeneous orbit.
\end{lemma}

\begin{proof}
At connected monodromy level decompose the complex representation
into $G_{c,\C}$-isotypical summands.  Each is stable under the Hodge
circle, since that circle acts on the connected group by inner
automorphisms.  An isotypical summand is the tensor product of an
irreducible representation with its multiplicity space.  After a
finite covering of the circle, lift its inner action to the first
factor; the remaining circle action is on the multiplicity space.
Its eigenspace decomposition permits a decomposition into
group-irreducible summands stable under the circle.  The corresponding
flat projectors are Hodge-compatible at a basepoint and, by
\Cref{lem:marked-monodromy}, everywhere.  At least one projector
has generically nonzero value on $u$.

Conjugate this projector by the full monodromy group.  Its orbit is
finite because $G_c$ acts trivially on the projector.  Tensor the
projected vectors over that finite orbit.  On the normal connected
level all factors are generically nonzero: the deck transformations
carry one projected vector to the others.  Permutation of the
factors makes the tensor construction equivariant for full
monodromy, so it descends.  Its source is a positive tensor power
of $\mathcal L$.

For the connected group take the Cartan component, the irreducible
summand of highest weight equal to the sum of the highest weights
of the factors.  It has multiplicity one, and its projection is
preserved by simultaneous automorphisms and permutations of the
factors.  This projection never kills a pure tensor of nonzero
vectors.  Indeed, for a nonzero vector in an irreducible
representation its coefficient along an extremal weight line
after translation by $g\in G_{c,\C}$ is a nonzero regular function
of $g$.  Finitely many such functions are simultaneously nonzero
on a nonempty open set.  Their product is the corresponding
extremal coefficient in the Cartan component.  Since the
projection commutes with translation, its value on the original
tensor was nonzero as well.  The projector preserves the Hodge
grading.  The projected vector is in the highest Hodge degree of
the tensor, which, because its projection is nonzero, is also the
highest degree of the retained Cartan component.

Retain each rational simple adjoint factor having at least one
complex simple ideal acting nontrivially on this irreducible
representation.  This collection is full-monodromy invariant.
Let $q:M\to\operatorname{Aut}(\mathfrak g)$ be the projected
adjoint action, and put $K=\ker q$.  Every element of $K$
commutes with the represented Lie algebra.  Schur's lemma implies
that its action on the irreducible representation is scalar.
Every ideal in $\Lie K$ is unrepresented by the definition of the
retained collection, so $K^\circ$ acts trivially.  Consequently
the image of $K$, although $K$ itself need not be finite, is a
finite group of scalars.  Choose a positive integer killing that
image and take the Cartan component in the corresponding tensor
power.  The resulting representation factors through $q(M)$,
and the preceding nonvanishing argument applies again.  In
particular its connected representation factors through the
adjoint group $G$.

These constructions, as well as subsequent invariant projections,
are homogeneous of positive degree in $u$.  The number of retained
rational factors is a nonnegative integer, so it has a minimum
among these choices.  Fix a minimizing choice.  Its purpose is to
prevent a further nonzero invariant projection from discarding a
factor.  When a factor fixes the highest space throughout its orbit,
this minimality will force the instability that supplies an additional
flag.

Connected irreducibility does not require the highest Hodge step
$E_0$ to have rank one; the construction keeps the entire step.
The adjoint representation is faithful on $G$.  Tensor generation
for a reductive group, \cite[Theorem~4.14]{Milne}, realizes each
of its representations as a summand of finite sums of mixed
tensors of this faithful representation.  For a factor of $G$
the Hodge grading on its Lie algebra is a cocharacter of the
adjoint group.  Its action in an irreducible representation and
the grading inherited from $\mathbb A$ induce the same action on
the Lie algebra; their quotient is scalar.  They thus differ
only by a constant Hodge shift.  Equivariant projectors on the
Lie tensors preserve the grading, and the induced compact-form
metrics have the asserted compatibility with adjoints.  Applying
this separately to the factors gives the tensor decomposition of
$E_0$.  Finally the parabolic associated with the Hodge grading
preserves its maximal-weight subspace.  The stabilizer of $E_0$
therefore contains a parabolic, proving that its orbit is closed
and projective.

All projectors used above are flat Hodge morphisms.  They and all
tensor operations extend on unipotent canonical extensions, by
the functoriality recalled in \Cref{sec:preliminaries}.  The
Leibniz rule for the Higgs field proves the last assertion of
\Cref{prop:marked-period} about iterates.
\end{proof}

The homogeneous replacement has removed the scalar kernel while retaining
a nonzero polynomial expression in the original vector. Full periods
can still miss a direction of that vector on a compact factor. We now
record those directions as additional algebraic flags. They will allow
the joint quotient to retain the vector directions missed by the full
period map; a point quotient will give a constant coefficient system.

\begin{remark}\label{rem:prelim-algebraic-monodromy-torsor}
We will use an algebraic monodromy torsor both for the compact-flag
construction below and for the later horizontal-germ estimate
\Cref{lem:connection-dimension}. Here is its construction.
Restrict a variation to a smooth algebraic base and pass to a finite
étale cover killing the component group of its algebraic monodromy.
Suppose its connected algebraic monodromy $H\subseteq\operatorname{GL}(V)$
is semisimple.  By the stabilizer-of-a-line theorem for algebraic groups
and the tensor realization of representations of $\operatorname{GL}(V)$
\cite[Theorems 4.27 and 4.14]{Milne}, there is a finite tensor construction containing
a line whose stabilizer is exactly $H$.  Since a connected semisimple
group has no nontrivial character, $H$ fixes a nonzero vector $t$ in
that line, and the stabilizer of $t$ is still exactly $H$.  This vector
defines a global monodromy-invariant flat tensor section.  It is
algebraic by the full faithfulness of the regular-singular
Riemann--Hilbert correspondence \cite[Chapter II, \S6]{DeligneRSE}.
Frames carrying $t$ to this section form an algebraic $H$-torsor inside
the frame bundle; on a dense open it is a principal bundle, and its
connection is algebraic and preserves the reduction.  Its monodromy is
Zariski dense in $H$ by definition.  Thus regular singularity is used
to obtain the algebraic reduction, while the dimension estimate itself
only needs the flat algebraic connection in its statement.
\end{remark}


\subsection{Compact flags attached to the vector}

Call a retained rational simple factor \emph{invisible} if all its
complex factors act trivially on the orbit of $E_0$.  This condition
is invariant under the component group of the full monodromy.

\begin{lemma}\label{lem:marked-flags}
For the minimizing replacement of \Cref{lem:marked-replacement}
there is a generically algebraic section $j_{\mathrm u}$ with the
compact-flag and factor-visibility properties in
part~\textup{(ii)} of \Cref{prop:marked-period}.
Its factors have flat positive definite metrics, and their flag
varieties carry positive Plücker lines with invariant metrics.
The construction is pointwise equivariant in $[u_*]$.
\end{lemma}

\begin{proof}
On a nontrivially represented complex ideal in an invisible rational
factor the maximal Hodge subspace is invariant under the whole
simple group.  Irreducibility makes it the whole representation
factor.  The grading on that representation is scalar, and
therefore its induced Lie grading is zero.  The Hodge metric on
this Lie ideal is consequently flat and positive definite.
Monodromy acts isometrically for this compact structure.

Fix an orbit of invisible rational factors under the full
monodromy, and let $U$ be the complexification of their product.
Consider the action of $U$ on a general nonzero value of $u_*$.
If this vector were not unstable for the origin, some homogeneous
invariant polynomial of positive degree would be nonzero on it.
Equivalently, its tensor power would have a nonzero projection to
the $U$-invariant subspace.  This elementary invariant-theoretic
equivalence follows because invariants separate the disjoint
closed sets $\{0\}$ and a closed orbit in the orbit closure;
averaging over a maximal compact subgroup gives the invariant
projection.  The projection is Hodge-compatible: $U$ is normalized
by the Hodge circle.  It is also compatible with full monodromy,
because the chosen product of rational factors is invariant under
it.  After repeating \Cref{lem:marked-replacement}, the nonzero
invariant projection eliminates all factors in this orbit and
introduces none.  This contradicts minimality.  Thus the vector
is generically unstable for each such $U$.

Choose on the cocharacter lattice of $U$ an integral positive
length form invariant under its Weyl group and under the finite
factor automorphisms in question.  One obtains it by summing
positive invariant forms over those automorphisms.  Kempf's
optimal one-parameter subgroup theorem
\cite[\S3]{Kempf78}, in its origin-instability form
\cite[Theorem~6]{Kempf76}, applied to the affine representation
and the closed origin, assigns to an unstable vector a unique
optimal parabolic.  It is equivariant under $U$ and under every
automorphism preserving the representation and the length form.
It is unchanged by multiplying the vector by a nonzero scalar.
It is proper, since $U$ is semisimple and an optimal positive
slope requires a noncentral cocharacter.

We spell out why this pointwise assignment supplies an algebraic
section on a dense open.  Fix a maximal torus.  There are finitely
many subsets of the weight set of the representation.  For each
subset, maximizing the positive normalized minimum pairing is
the closest-point problem for its convex hull, with respect to
the chosen rational positive form.  It determines a rational
optimal ray when the minimum is positive.  Thus there are only
finitely many optimal cocharacter directions and slopes up to
conjugacy.  The conditions that a translate of a vector have a
specified nonzero weight support are constructible.  Their
images, and the conditions selecting the maximal slope, are
constructible by Chevalley's theorem.  The resulting graph of
optimal parabolics on the unstable locus is constructible and single
valued.  Use the algebraic principal monodromy reduction of
\Cref{rem:prelim-algebraic-monodromy-torsor}; its associated projective
representation and flag bundles are algebraic.  Equivariance makes the
graph an algebraic constructible subset of their product over the base.
Pull this graph back along the actual algebraic section $[u_*]$, on the
dense open where that section is unstable.  The resulting constructible
graph has exactly one point over each point of this open base.  Over an
irreducible base in characteristic zero, its dominant component is
birational to the base: after restricting to dense opens, its projection
is a generically injective morphism and therefore has degree one.
It consequently defines a rational section, regular after shrinking.
Applying this argument to the actual section avoids any assumption that
it misses the indeterminacy locus of a rational map on the ambient
projective representation bundle.  The construction uses algebraic associated
bundles, not algebraic local trivializations by holomorphic flat frames.

An optimal cocharacter has zero component on every ineffective
factor, since deleting such a component preserves its slope
numerator and decreases its length.  All nontrivial flags
therefore occur on the isometric complex ideals identified
above.  On the irreducible connected level the generic parabolic
types are fixed.  The component action permutes these types
along with the rational factors.  Properness on at least one
factor in an orbit implies properness on some complex factor of
each rational factor in that orbit.  Repeating this for all
invisible orbits gives $j_{\mathrm u}$.

The highest-space marks already see every other rational factor.
Stabilizers of highest spaces and of the new flags are parabolic,
and the two constructions use disjoint factors.  Thus their
joint orbit is the asserted product $J$.  A tensor product of
highest subspaces determines each factor subspace, so this
description is also valid when $E_0$ is presented as a single
subspace of $\mathbb A$.  Embed a parabolic flag by the top
exterior power of its parabolic subalgebra.  The dual tautological
Plücker line is positive, with the metric induced by the flat
compact metric.  Tensoring these lines gives a positive line
invariant also under the relevant factor permutations.
\end{proof}

The marks now see every retained rational factor. We next construct an
algebraic parameter space and descend the actual coefficient variation
to it. Descent of the filtration alone would leave a possible finite
scalar obstruction, which is why that kernel was removed first.

\Needspace{10\baselineskip}
\subsection{An algebraic quotient and actual descent}

\begin{lemma}[Generic Lie stabilizer]\label{lem:generic-lie-stabilizer}
Let a polarized rational Lie variation on a smooth connected complex
quasi-projective variety have an invariant integral lattice, semisimple
fiber $\mathfrak a$, and connected algebraic monodromy
$\operatorname{Int}(\mathfrak a)$.  At a Hodge-generic lifted
period, a rational Lie automorphism preserving the Hodge
filtration is the identity.
\end{lemma}

\begin{proof}
A rational endomorphism preserving the filtration of a pure Hodge
structure also preserves its conjugate filtration, hence is a
Hodge endomorphism.  It commutes with the Mumford--Tate group.
At a Hodge-generic point this is the generic Mumford--Tate group,
which contains the connected monodromy by
\cite[\S5, Theorem~1]{Andre92}.  If $\alpha$ is also a Lie
automorphism, then for every $x\in\mathfrak a$,
\[
 \operatorname{ad}(\alpha x)
   =\alpha\operatorname{ad}(x)\alpha^{-1}
   =\operatorname{ad}(x).
\]
The Lie algebra has zero center, so $\alpha x=x$.  Such points
may be chosen while imposing finitely many generic differential
rank conditions: in a flat chart the proper loci where any one
of the countably many rational tensors acquires an additional
Hodge condition can be avoided together with those conditions.
\end{proof}

\begin{lemma}[Period quotient and descent]\label{lem:period-descent}
There is an algebraic quotient with connected generic fibers
which remembers the full integral period of $\mathfrak g$ and
the compact marks.  Both $\mathfrak g$ and $\mathbb A$ descend
as variations on a dense smooth open of the quotient.  Its full
Lie period and compact marks give a generically immersive joint
map in flat coordinates.  The same construction is available for
any monodromy-stable collection of rational normal factors, at
connected level if necessary.
\end{lemma}

\begin{proof}
Use the full filtration on the rational integral polarized Lie
variation, not only the highest subspace of $\mathbb A$.  Let
$\Gamma$ be an arithmetic group of isometries of its polarized
lattice containing its monodromy, and let $\Omega$ be its period
domain, or the associated generic Mumford--Tate domain.
The algebraicity theorem for pure integral period maps
\cite[Theorem~1.1]{BBT} factors the period map through a
dominant algebraic map to a quasi-projective algebraic image.
Its hypotheses hold here by \Cref{lem:marked-monodromy}; in
particular a lattice has been retained in the full Lie system.
Normalize this image in the relative algebraic closure of its
function field in $\C(B)$.  The resulting map
$b_0:B^\circ\dashrightarrow T^\circ$ has geometrically connected
generic fiber.  Shrink to smooth opens where it is a morphism.
The dimension of $T^\circ$ is the generic rank of the full
period map.  Indeed, the real period stabilizer is compact,
arithmetic orbits are discrete, and the factorization preserves
the local analytic image dimension.

Here is the local-system descent, including possible finite
monodromy.  Compactify the graph properly over $T^\circ$, resolve
the complement, and shrink $T^\circ$ so that the compactification
and all strata of its SNC complement are smooth over it.
Lifting vector fields tangent to the strata gives a local
differentiable trivialization of this proper smooth pair and
hence of its open part.  We may thus assume that $b_0$ is a
topological fiber bundle with connected fibers.  Removing a
proper algebraic subset of a smooth complex variety induces a
surjection on its fundamental group, so these shrinkings do not
decrease the relevant monodromy images.

Choose a Hodge-generic point of $B^\circ$.  Along the universal
cover of its $b_0$-fiber the lifted Lie-period map has values in
one discrete arithmetic orbit; it is therefore constant.
Every vertical loop acts by a rational Lie automorphism fixing
that lifted period.  \Cref{lem:generic-lie-stabilizer} makes its
action trivial.  The exact sequence
\[
 \pi_1\bigl(b_0^{-1}(t)\bigr)\longrightarrow
 \pi_1(B^\circ)\longrightarrow\pi_1(T^\circ)
 \longrightarrow1
\]
now descends the Lie representation.  The representation on
$\mathbb A$ descends as well: by
\Cref{lem:marked-replacement} it factors through the full
projected Lie action, including its component group.  This is
where removal of the finite scalar image is necessary.

The Lie filtration is constant in flat coordinates along each
connected fiber.  The derivative of the filtration on
$\mathbb A$ is the action of
\eqref{eq:marked-period-tangent}; it too vanishes on those
fibers.  Consequently both filtrations descend.  Holomorphicity,
Griffiths transversality, and polarization can be checked after
the smooth surjective pullback.  The descended local systems
have quasi-unipotent local monodromy: for a boundary valuation
of $T$ choose a valuation above it on a proper model of $B$;
a positive power of its monodromy is a monodromy of the original
variation.  Quasi-unipotence of that power implies
quasi-unipotence of the original operator.  Canonical extension
and GAGA, or regular-singular Riemann--Hilbert, therefore give
the algebraic flat bundles and algebraic Hodge bundles on these
opens; see \cite[Chapter~II]{DeligneRSE} and
\Cref{sec:preliminaries}.

Over $T^\circ$ consider the algebraic flag bundle of the
descended compact factors.  The section $j_{\mathrm u}$ gives
an algebraic map from $B^\circ$ into that bundle.  Take its
image and again normalize in the relative algebraic closure
inside $\C(B)$.  This is the desired $S^\circ$.  Variations
pull back from $T^\circ$, and the compact flag descends by its
tautological definition.  On a dense open the finite map from
$S^\circ$ to its image is étale, and the period-image
factorization preserves tangent rank.  Hence the full period
and compact marks have joint differential of rank $\dim S$.
Resolve projective compactifications and the map from $B$ to
obtain $b:B\to S$ with connected fibers.  The monodromy images
on the open source and target agree by the fundamental-group
sequence just used.

For a monodromy-stable collection of rational factors the
projected Lie system again has its rational polarization and
lattice.  The identical argument gives its quotient and
descent.  At connected level every rational normal factor
collection is permitted.
\end{proof}

\phantomsection\label{proof:point-period}
If $S$ is a point, the full projected period is constant.
Its monodromy lies in a compact period stabilizer and in the
discrete integral isometry group, and is thus finite.  Its
connected algebraic monodromy $G$ is trivial.  The representation
on $\mathbb A$ factors through the projected Lie action, whose
kernel has already been removed, so $\mathbb A$ is constant.
When there are no retained factors the construction simply takes
$S$ to be a point.

\subsection{The metric comparison used for nefness}

The quotient has now been constructed and the coefficient variation
descends to it.  Two different differentials enter the positivity
argument.  The full joint differential $(t,dj_{\mathrm u})$ has generic
rank $\dim S$, by \Cref{lem:period-descent}; it will control the
logarithmic adjoint.  The differential of the marks
$j=(E_0,j_{\mathrm u})$ can have smaller rank.  It is this second
differential that the line $\det E_0+\mathcal H_{\mathrm u}$ measures
on the open set.  We first compute its curvature and then use the
boundary growth estimates to identify curvature rank with numerical
dimension on a projective model.

In the smooth Hodge splitting $E=\bigoplus_p E^p$, write
\[
 \nabla=\mathcal D+\theta+\theta^\dagger,
\]
where $\mathcal D$ is the graded Chern connection and
$\theta^\dagger$ is the Hodge adjoint of the degree-lowering Higgs field.
If $\theta_p(v):E^p\to E^{p-1}$, flatness and the polarization give,
with positive curvature normalization,
\begin{equation}\label{eq:prelim-hodge-curvature}
 R_{E^p}(v,\bar v)=
 \theta_p(v)^*\theta_p(v)
 -\theta_{p+1}(v)\theta_{p+1}(v)^*.
\end{equation}
The irrelevant common positive normalization factor is suppressed here.
For the highest nonzero step $F^p=E^p$, the second term vanishes, and hence
\begin{equation}\label{eq:prelim-top-curvature}
 c_1(\det F^p,h)(v,\bar v)
 =c\,\|\theta_p(v)\|_{\mathrm{HS}}^2\quad(c>0).
\end{equation}
In a local flat frame the differential of the subspace map
$b\mapsto F^p_b$ is exactly $\theta_p$; transversality and $F^{p+1}=0$
ensure that its image lies in $F^{p-1}/F^p$.  Thus the curvature kernel
in \eqref{eq:prelim-top-curvature} is exactly the kernel of that
subspace differential.  This statement concerns the whole highest
step, not a chosen line inside it.  Adding the positive compact-flag
form gives kernel $\ker(dE_0,dj_{\mathrm u})$.

To extend this computation across the boundary, we use the
several-variable norm estimates of \cite[Theorem 5.21]{CKS}, also stated
and proved in \cite[Theorem 5.1]{CK}.  Work at unipotent level with the
canonical Hodge extensions of \Cref{sec:preliminaries}.  In a boundary
chart $\Delta=\{z_1\cdots z_\ell=0\}$, shrink so that $|z_i|<e^{-1}$
and put $\Lambda=\prod_{i=1}^{\ell}(-\log|z_i|)$.  Relative to any
smooth positive metric $h_0$ on the extending bundle there are constants
$C,N>0$ such that its Hodge metric satisfies
\begin{equation}\label{eq:prelim-hodge-growth}
 C^{-1}\Lambda^{-N}h_0\ \leq\ h\ \leq\ C\Lambda^N h_0.
\end{equation}
The same assertion holds for the dual, every extending filtration step,
and every Hodge graded quotient, with their induced metrics.  To pass
from the usual sector estimates to \eqref{eq:prelim-hodge-growth}, divide
the logarithmic variables into their finitely many order sectors.  Each
ratio power in the sector estimate is bounded by a power of $\Lambda$.
The finite logarithmic exponential that converts flat multivalued frames
to canonical frames has entries polynomial in the logarithms.  Applying
the same bounds to the dual gives the lower estimate.  Restriction to a
subbundle and passage to its quotient preserve a two-sided comparison
with a smooth metric, which gives the assertions for $F^p$ and $\Gr_F^p$.
The extension compatibility with flat Hodge projectors proved in
\Cref{sec:preliminaries} gives the same estimates for the complex Hodge
summands used here.

The induced line metrics therefore have two-sided logarithmic-logarithmic
weight bounds.  A curvature computation on the open set alone would not
identify their intersection numbers.  The following full-mass argument
makes that comparison.


\begin{lemma}\label{lem:log-metric}
Let $Z$ be a smooth projective variety of dimension $n$, let $\Delta$ be
an SNC divisor, and let $L$ be a rational line bundle.  Suppose that on
$Z^\circ=Z\setminus\Delta$ it has a smooth Hermitian metric with
semipositive curvature form $\Theta$, normalized to represent $c_1(L)$.
Suppose that in regular nonvanishing local frames of a Cartier multiple
of $L$ the metric weights satisfy
\begin{equation}\label{eq:prelim-loglog-bound}
 |\varphi(z)|\leq C\left(1+
       \sum_{i=1}^{\ell}\log(-\log|z_i|)\right)
\end{equation}
near every boundary chart $\Delta=\{z_1\cdots z_\ell=0\}$.
Then $L$ is nef.  For every Kähler form $\eta$ on $Z$ and
$0\leq k\leq n$,
\begin{equation}\label{eq:prelim-mass-identity}
 \int_{Z^\circ}\Theta^k\wedge\eta^{n-k}
       =c_1(L)^k[\eta]^{n-k}.
\end{equation}
In particular,
\begin{equation}\label{eq:prelim-curvature-rank}
 \nu(L)=\max_{z\in Z^\circ}\rk\Theta_z.
\end{equation}
When the curvature is real analytic, as for the induced Hodge and
subspace metrics used below, this maximum is its generic rank.
\end{lemma}

\begin{proof}
Taking a Cartier multiple and dividing its weights and curvature by that
multiple reduces to a line bundle.  We use $dd^c$ normalized so that the
curvature of a metric of weight $\varphi$ is $dd^c\varphi$.

First we extend its local plurisubharmonic weights.  For $\epsilon>0$,
the function
\[
 \varphi_\epsilon=\varphi+
               \epsilon\sum_{i=1}^{\ell}\log|z_i|
\]
is plurisubharmonic off the coordinate divisor and is locally bounded
above near it, by \eqref{eq:prelim-loglog-bound}.  The removable
singularity theorem extends it plurisubharmonically to the polydisk.
On a smaller polydisk, the maximum principle, applied successively in
the coordinate variables, bounds it above by its values on a fixed
distinguished boundary torus disjoint from $\Delta$.  These bounds are
uniform as $\epsilon\downarrow0$.  Its increasing limit, with upper
semicontinuous regularization, is therefore a plurisubharmonic extension
of $\varphi$.  The extensions agree under the pluriharmonic frame changes
and give a closed positive current $T\in c_1(L)$ on $Z$.

This current has zero Lelong number at every point.  At a boundary point
approach along a path on which all the vanishing boundary coordinates
have absolute values comparable to a common parameter $r\downarrow0$.
The two-sided bound in \eqref{eq:prelim-loglog-bound} is then
$O(\log\log(1/r))=o(|\log r|)$.  The growth characterization of the
Lelong number gives an upper bound of zero; positivity gives the reverse
bound.  Away from $\Delta$ the current is smooth.  Demailly's
regularization theorem, in its zero-Lelong-number consequence
\cite[Corollary 6.4]{Demailly92}, now shows that $c_1(L)$ is nef.

We next prove the mass assertion, since zero Lelong numbers alone would
not suffice for it.  Fix $\epsilon>0$.  Nefness makes
$\alpha_\epsilon=c_1(L)+\epsilon[\eta]$ a Kähler class.  Choose a Kähler
representative $\omega_\epsilon$ and write
\[
 T+\epsilon\eta=\omega_\epsilon+dd^c\phi.
\]
Here $\phi$ is a global $\omega_\epsilon$-plurisubharmonic function,
smooth on $Z^\circ$; changing the smooth background has not changed its
logarithmic growth bound.  Write the irreducible components of $\Delta$
as $\Delta_i$, choose their defining sections $s_i$ with smooth metrics,
and rescale these metrics so that
\[
 u_i=-\log\|s_i\|^2\geq1\quad\text{on }Z^\circ.
\]
For any fixed $0<b<1$ and sufficiently small $\delta>0$, the function
\begin{equation}\label{eq:prelim-mass-barrier}
 \psi=-\delta\sum_i u_i^b
\end{equation}
is $\omega_\epsilon$-plurisubharmonic.  Indeed, on $Z^\circ$ each
$dd^c u_i$ is a fixed smooth form and
\[
 dd^c(-\delta u_i^b)
   =-\delta b u_i^{b-1}dd^c u_i
     +\delta b(1-b)u_i^{b-2}du_i\wedge d^c u_i.
\]
The second summand is positive and the first is bounded below by
$-\delta C\omega_\epsilon$, because $u_i\geq1$.  Decrease $\delta$ so
that the sum of these lower bounds is at least
$-\tfrac12\omega_\epsilon$.  The same extension argument applies across
$\Delta$.  Since $\log u=o(u^b)$, the lower growth bound for $\phi$
also gives a constant $C_\epsilon$ such that
\begin{equation}\label{eq:prelim-less-singular}
 \phi\geq\psi-C_\epsilon.
\end{equation}

We verify directly that $\psi$ has full Monge--Ampère mass.  Take a smooth
decreasing function $\chi:[0,\infty)\to[0,1]$ that is $1$ on $[0,1]$
and $0$ on $[2,\infty)$.  For $R>1$ set
\[
 f_R(t)=1+\int_1^t b s^{b-1}\chi(s/R)\,ds\qquad(t\geq1).
\]
Then $f_R=t^b$ for $t\leq R$, $f_R$ is increasing and concave, and it
is constant for $t\geq2R$.  Its derivative bounds, uniformly in $R$, are
\begin{equation}\label{eq:prelim-barrier-truncation}
 0\leq f_R'(t)\leq C t^{b-1},\qquad
 0\leq-f_R''(t)\leq C t^{b-2}.
\end{equation}
Moreover $f_R\uparrow t^b$ as $R\to\infty$.  The functions
$\psi_R=-\delta\sum_i f_R(u_i)$ extend smoothly over $Z$, are
$\omega_\epsilon$-plurisubharmonic for the same choice of $\delta$,
and equal $\psi$ locally on $Z^\circ$ once $R$ is sufficiently large.

In a fixed boundary coordinate chart, comparison of
$du_i=-d\log|z_i|^2+$ a smooth form and
\eqref{eq:prelim-barrier-truncation} gives a bound of positive forms
\begin{equation}\label{eq:prelim-local-majorant}
 0\leq\omega_\epsilon+dd^c\psi_R
 \leq C\left(\eta_0+
       \sum_{i=1}^{\ell}
       \frac{\sqrt{-1}\,dz_i\wedge d\bar z_i}
       {|z_i|^2(-\log|z_i|)^{2-b}}\right),
\end{equation}
where $\eta_0$ is a smooth positive coordinate form and $C$ is independent
of $R$.  The right-hand side has an integrable $n$th power.  Each
singular summand has rank one and its square is zero, and its normal
integral is bounded by a constant times
\[
 \int_0^{r_0}\frac{dr}{r(-\log r)^{2-b}}<\infty.
\]
Products involving different normal coordinates are integrable by
Fubini.  A finite chart covering and dominated convergence now imply
\begin{align*}
 \int_{Z^\circ}(\omega_\epsilon+dd^c\psi)^n
  &=\lim_{R\to\infty}
        \int_Z(\omega_\epsilon+dd^c\psi_R)^n\\
  &=\int_Z\omega_\epsilon^n.
\end{align*}
The last equality is Stokes' theorem for the smooth approximants.  Since
the non-pluripolar Monge--Ampère product does not charge $\Delta$, this
is exactly full mass for $\psi$.

The full-mass class is preserved on passing to a less singular
plurisubharmonic potential.  Apply
\cite[Proposition 1.6]{GZ} to the Kähler form $\omega_\epsilon$ and
\eqref{eq:prelim-less-singular}.  All the hypotheses are satisfied:
$Z$ is compact Kähler, both potentials are
$\omega_\epsilon$-plurisubharmonic, and the lower potential has just
been shown to have full mass.  It follows that
\begin{equation}\label{eq:prelim-epsilon-mass}
 \int_{Z^\circ}(\Theta+\epsilon\eta)^n
       =(c_1(L)+\epsilon[\eta])^n
       \qquad(\epsilon>0).
\end{equation}
All mixed integrands on the left are nonnegative.  The identity for one
positive $\epsilon$ therefore proves that each has finite integral.
Expansion of \eqref{eq:prelim-epsilon-mass} and comparison of polynomial
coefficients proves \eqref{eq:prelim-mass-identity}.

Finally, a smooth semipositive $(1,1)$-form has
$\Theta^k\wedge\eta^{n-k}>0$ at a point exactly when its rank there is
at least $k$.  If this happens at one point it happens on an open
neighborhood, so its integral is positive.  Combining this observation
with \eqref{eq:prelim-mass-identity} and
\eqref{eq:prelim-numerical-dimension} proves
\eqref{eq:prelim-curvature-rank}.  For a real-analytic form the nonzero
maximal minors cannot vanish on an open set, so the maximal-rank locus
is dense.  This proves the final assertion.  When
$n=0$ the statements follow from the empty-product convention.
\end{proof}

\begin{remark}\label{rem:prelim-metric-operations}
Estimate \eqref{eq:prelim-hodge-growth} supplies
\eqref{eq:prelim-loglog-bound} for determinants of extending Hodge
bundles, their holomorphic subbundles and quotients, and duals and tensor
products of these lines.  For a subspace given initially only on a dense
open set, first resolve its rational Grassmann map; the resulting
tautological subbundle is an actual subbundle of the pulled-back
extending bundle.  Its inclusion is uniformly nondegenerate relative
to smooth metrics, so the two-sided estimate applies.  Multiplication
by a smooth positive metric factor only adds a bounded weight.  The
semipositivity required by \Cref{lem:log-metric} must still be proved for
each line by its curvature calculation.

The rank computation may also be made on a smooth generically finite
cover $q:Z'\to Z$.  If the lemma applies to $q^*L$, it gives nefness of
$q^*L$, hence of $L$, by testing a curve using a curve lying over it.
To compute intersections downstairs, the mass identity upstairs remains
valid with $q^*\eta$ in place of the Kähler form: apply it first to
$q^*\eta+t\eta'$ for $t>0$, where $\eta'$ is Kähler, and compare
coefficients in $t$.  The projection formula gives
\[
 (q^*c_1(L))^k(q^*[\eta])^{n-k}
       =\deg(q)\,c_1(L)^k[\eta]^{n-k}.
\]
The form $q^*\eta$ is positive definite on the étale locus, which is
dense.  A nonempty open maximal-curvature-rank locus meets it.  Thus the
same curvature rank computes $\nu(L)$ downstairs.
\end{remark}


\subsection{The nef line and an initial logarithmic adjoint}

Assume $\dim S>0$.  Choose a normal finite monodromy level of
$S^\circ$ contained in the connected group and in a principal
integral congruence subgroup of sufficiently high level, for
example a level divisible by $3$.  Quasi-unipotent elements at
this level are unipotent.  One elementary way to see the
assertion is that if $A\equiv1\pmod m$, then
$(\zeta-1)/m$ is an algebraic integer for each root-of-unity
eigenvalue $\zeta$.  For $m\ge3$ all its conjugates have
absolute value less than one unless it is zero; hence
$\zeta=1$.

Take an equivariant smooth projective compactification
$\widetilde S$ of this level, with SNC boundary.  On the compact
complex ideals the now-unipotent local monodromy is unitary and
therefore trivial.  Their flat bundles extend across the
boundary with their smooth positive metrics.  Resolve the
compact-flag maps in these bundles equivariantly.  The line
\begin{equation}\label{eq:marked-level-line}
 \widetilde{\mathcal H}
   =\det\widetilde E_0+\widetilde{\mathcal H}_{\mathrm u}
\end{equation}
is defined on this model: $\widetilde E_0$ is the highest
Schmid extension, and
$\widetilde{\mathcal H}_{\mathrm u}$ is the product Plücker
line from \Cref{lem:marked-flags}.  Both have semipositive
curvature.  For $\det\widetilde E_0$ the curvature is the
squared norm of the differential of the highest subspace; for
the compact summand it is the pullback of the positive flag
form.  Their kernels therefore intersect in the kernel of $dj$.
The Hodge norm and dual norm have logarithmic growth in
canonical-extension frames; the compact metric is smooth.
\Cref{lem:log-metric} gives nefness and
\begin{equation}\label{eq:marked-nef-rank}
 \nu(\widetilde{\mathcal H})
   =\max\operatorname{rank}(dj).
\end{equation}

The construction is linearized for the deck group.  A common
positive power kills the finite stabilizer actions on the
fibers of each line, and that power descends to the finite
quotient.  For clarity, descent can be checked locally by
averaging a local frame after its stabilizer acts trivially;
the invariant frame is nonzero after shrinking and supplies
the required quotient frame.  Let $S$ now be a smooth
projective log resolution of that quotient, replacing the
earlier compactification, and pull the descended rational
lines to $S$.  Denote them by $\mathcal H$ and
$\mathcal H_{\mathrm u}$.  Nefness descends under a finite
surjection and is preserved by pullback.  On the resolution
of the corresponding base change upstairs the lines are
exactly the pullbacks of \eqref{eq:marked-level-line}, by
functoriality of unipotent extensions and of the resolved flag
maps.  In particular
\begin{equation}\label{eq:marked-downstairs-rank}
 \nu(\mathcal H)=\max\operatorname{rank}(dj),
 \qquad
 \mathcal H_{\mathrm u}\text{ and }
 \mathcal H-\mathcal H_{\mathrm u}\text{ are nef}.
\end{equation}

We record compatibility with rational normal factors, used in
the next section.  At connected level, write $G=H\times Q$.
Then $\mathbb A=A_H\otimes A_Q$ and
$E_0=E_H\otimes E_Q$.  Splitting the compact marks accordingly
gives
\begin{align}\label{eq:marked-factor-lines}
 \mathcal H_H
   &=(\rk E_Q)\det E_H+\mathcal H_{\mathrm u,H},&
 \mathcal H_Q
   &=(\rk E_H)\det E_Q+\mathcal H_{\mathrm u,Q},&
 \mathcal H&=\mathcal H_H+\mathcal H_Q.
\end{align}
These are nef lines on suitable smooth models.  Apply
\Cref{lem:period-descent} to the full $Q$-period and its
compact marks, obtaining $q:S\to S_Q$ after modification.
The $Q$-line is $q^*\underline{\mathcal H}_Q$ for the same
construction on $S_Q$.  The restricted congruence condition
holds on the lattice intersected with each rational ideal, so
the descended $Q$-system has unipotent boundary monodromy.
Consequently its extensions, flag lines, and these identities
remain compatible on a common resolution.

Mark the reduced SNC union $D_0$ of those boundary prime
divisors of $S$ whose local monodromy on $\mathfrak g$ is
not the identity.

\begin{lemma}[Initial adjoint bigness]\label{lem:initial-adjoint}
For the marked-period quotient just constructed,
\[
 K_S+D_0+c\mathcal H_{\mathrm u}
 \quad\text{is big for all sufficiently large }c.
\]
\end{lemma}

\begin{proof}
Let $d=\dim S$ and work first on a smooth unipotent level
$p:\widetilde S\to S$, resolving further when necessary.
On its open part the joint differential has image
\[
 \mathcal J\subseteq
 \Gr_F^{-1}\mathfrak g_{\C}\oplus
 j_{\mathrm u}^*T_{J_{\mathrm u}}.
\]
The differential here is $(t,dj_{\mathrm u})$, where the
second component is the vertical derivative in flat flag
charts.  It is globally defined because the transition maps
are flat.  Its rank is $d$ by \Cref{lem:period-descent}.
Resolve the map to the corresponding relative Grassmannian
so that $\mathcal J$ extends as a subbundle of the indicated
extended ambient bundle.

Equip it with the subbundle metric.  We give the curvature
calculation because it uses the full period differential,
whereas the nef line $\mathcal H$ only uses the marks $j$.
Fix a tangent vector $v$, and choose an orthonormal frame
$(a_\mu,b_\mu)$ of $\mathcal J$ at the point.  Each
$a_\mu$ lies in the image of $t$, so integrability gives
$[t(v),a_\mu]=0$.  The Hodge curvature formula on the
degree $-1$ Lie space consequently gives
\[
 \sum_\mu
 \bigl\langle\Theta_{\Gr_F^{-1}\mathfrak g}(v,\bar v)
          a_\mu,a_\mu\bigr\rangle
 =-\sum_\mu\bigl\|[t(v)^*,a_\mu]\bigr\|^2.
\]
The convention is the curvature normalization of
\eqref{eq:prelim-hodge-curvature}.  The flag tangent curvature is
bounded above by $C\omega_{\mathrm u}(v,\bar v)$ times
the identity, where $\omega_{\mathrm u}$ is the positive
compact-flag form pulled back along the section.  Such a
uniform $C$ exists because the flag variety is compact and
the flat transition actions are isometric.  Subbundle
curvature subtracts a squared second fundamental form.
After increasing $C$ by the fixed rank $d$, this proves
\begin{equation}\label{eq:marked-curvature-estimate}
 \Theta_{-\det\mathcal J+c\widetilde{\mathcal H}_{\mathrm u}}
       (v,\bar v)
 \geq \sum_\mu\bigl\|[t(v)^*,a_\mu]\bigr\|^2
          +(c-C)\omega_{\mathrm u}(v,\bar v).
\end{equation}
The right side is positive whenever the joint differential
is nonzero.  Indeed, if its compact component is nonzero,
the last term is positive for $c>C$.  Otherwise
$t(v)\ne0$.  Since $(t(v),dj_{\mathrm u}(v))$ belongs
to $\mathcal J$, vanishing of all the commutators in
the first term would imply $[t(v)^*,t(v)]=0$.
But $t(v)$ acts as a nonzero lowering nilpotent in the
faithful adjoint representation.  A normal nilpotent
endomorphism is zero, a contradiction.  Thus the line
\[
 R=-\det\mathcal J+c\widetilde{\mathcal H}_{\mathrm u}
\]
has semipositive curvature and full generic curvature rank.
The extended ambient Hodge metrics and their duals have
logarithmic bounds, the compact metric is smooth, and the
resolved subbundle metric has the same bounds.  By
\Cref{lem:log-metric}, $R$ is nef with numerical dimension
$d$, hence big.

It remains to compare $R$ with the claimed adjoint on $S$.
The differential gives a generically invertible map
\begin{equation}\label{eq:marked-log-tangent-map}
 p^*T_S(-\log D_0)\dashrightarrow\mathcal J.
\end{equation}
It is regular at every prime upstairs which dominates a
prime of $S$.  At a marked prime, the finite cover is
generically a power substitution $t=z^e$.  The lift of
$t\partial_t$ is $e^{-1}z\partial_z$.
The logarithmic Higgs field extends on the unipotent
canonical extension, so its Lie component is regular;
one recovers the Lie-valued $t$ from the adjoint Higgs
field by the split adjoint inclusion into endomorphisms.
The resolved compact flag is regular upstairs, so its
derivative on $z\partial_z$ is regular as well.  The
tangential generators cause no poles.  At an unmarked
prime the monodromy level is generically unramified,
since it was defined using only the projected adjoint
monodromy.  The local Lie monodromy is the identity and
the extended period has no logarithmic term.  The same
regularity assertion follows using the ordinary tangent
generator there.

These regular ambient maps land in $\mathcal J$ at the
primes in question.  Indeed, $\mathcal J$ is a subbundle,
so its quotient in the ambient bundle is torsion-free;
a regular ambient section whose generic value lies in
$\mathcal J$ has zero image in that quotient.  Taking
determinants in \eqref{eq:marked-log-tangent-map} and
dualizing gives a generically nonzero map
\begin{equation}\label{eq:marked-adjoint-comparison}
 R\dashrightarrow p^*(K_S+D_0+c\mathcal H_{\mathrm u})
\end{equation}
regular above every prime of $S$.

Possible poles on divisors exceptional over $S$ do not
affect the conclusion, as follows.  Factor $p$ through
the finite normalization $S'\to S$ of this level.  Choose
an ample line $A_S$ on $S$ and a positive rational
$\epsilon$ sufficiently small that
$R-\epsilon p^*A_S$ is big.  For a sufficiently divisible
positive integer $m$, it has a nonzero section.  Its image
under \eqref{eq:marked-adjoint-comparison} is a rational
section of
\[
 p^*m(K_S+D_0+c\mathcal H_{\mathrm u}-\epsilon A_S)
\]
regular at every divisor of the normal variety $S'$.
Normality extends it over $S'$.  Taking its field norm
to $S$ gives a nonzero section of a positive multiple
of $K_S+D_0+c\mathcal H_{\mathrm u}-\epsilon A_S$.
Adding the ample term proves bigness.  Since
$\mathcal H_{\mathrm u}$ is nef, the conclusion persists
when $c$ is increased.
\end{proof}

This proves all construction and descent assertions of
\Cref{prop:marked-period}, as well as its initial adjoint
bigness.  The next section removes $D_0$ and proves the
tail-rank identity without any hypothesis that $u_*$ span
the whole of $E_0$.
